\documentclass[10pt,a4paper]{amsart}
\usepackage[margin=19mm]{geometry}
\usepackage{amsmath,amssymb,mathtools}
\usepackage[scr=rsfs,cal=euler]{mathalfa}
\usepackage{xfrac}
\usepackage[unicode,hidelinks,bookmarks=false]{hyperref}
\newcommand{\ie}{i.e.\ }
\newcommand{\cf}{cf.\ }
\newcommand{\resp}{respectively\ }
\newcommand{\Subsection}{\subsection}
\providecommand{\nobreakdash}{\nobreak\hbox{-}}
\setlength{\emergencystretch}{2em}
\multlinegap0pt
\usepackage{bm}
\usepackage{bbm}
\usepackage{dsfont}
\usepackage{xspace}
\usepackage{cancel}
\usepackage{mathtools}
\usepackage{amsthm}
\makeatletter
\@ifundefined{Theorem}{\newtheorem{Theorem}{Theorem}}{}
\@ifundefined{Exa}{\newtheorem{Exa}[Theorem]{Example}}{}
\@ifundefined{Prop}{\newtheorem{Prop}[Theorem]{Proposition}}{}
\makeatother
\DeclareMathOperator{\id}{\mathrm{id}}
\newcommand{\Corr}{\mathrm{Corr}}
\newcommand{\G}{\mathbb{G}}
\newcommand{\ovot}{\bar{\otimes}}
\title[Morita theory for dynamical von Neumann algebras]{Morita theory for dynamical von Neumann algebras}
\author{Joeri De Ro}
\date{Source: arXiv:2410.17407v3 (CC BY 4.0)}
\begin{document}
\maketitle

\noindent\emph{Source and licence.} Morita theory for dynamical von Neumann algebras. Version arXiv:2410.17407v3 (CC BY 4.0), distributed under the Creative Commons Attribution 4.0 International licence, as recorded in the arXiv licence metadata for this exact version and shown on the abstract page https://arxiv.org/abs/2410.17407v3. Full rights notice: licenses/arxiv-2410.17407-CC-BY-4.0.txt.\par\smallskip

\noindent\textit{Editorial scope.} Counted unit: the Theorem whose source label is \texttt{char2} in the article (source0.tex lines 1029--1040, source label \texttt{char2}), delivered together with the article's own complete original proof of it (source0.tex lines 1041--1154, one \texttt{proof} environment of 114 lines). No mathematics is written, repaired or re-derived: statement and proof are the article's own text line for line. Required context retained uncounted: amplification (lines 870--872); char1 (lines 887--893); main example (lines 984--993). Every definition, display and label of those lines is retained unchanged. Omitted from this package and not redistributed: 1--869, 873--886, 894--983, 994--1028, 1155--1722. Nothing inside a retained excerpt is omitted or altered: every retained body is delivered exactly as the article's lines are written, so no omission receipt of any kind accompanies this package. Citations: the retained excerpts carry no citation command at all, so no bibliography entry of the article is transcribed and no reference is redistributed. Reference adaptation: none was needed and none was made; every reference inside the delivered unit resolves inside it, so no unresolved reference or citation is produced. Printed slips kept literal: no printed word, symbol, hypothesis or number of the counted statement or of its proof was altered. Layout and class: the article's own class is \texttt{amsart} with packages including tikz-cd, fullpage, dsfont, hyperref, parskip, amsrefs, pdfpages, multirow, stmaryrd, array, amssymb, amsfonts, amsxtra, amsmath; the wrapper is the lane's pinned \texttt{amsart} editorial wrapper at 10pt on A4, which restates the theorem-like environments the retained text needs and adds the article's own macro definitions that the retained text needs (\texttt{\textbackslash Corr}, \texttt{\textbackslash G}, \texttt{\textbackslash id}, \texttt{\textbackslash ovot}), with extra wrapper packages (bm, bbm, dsfont, xspace, cancel, mathtools). The delivered numbering is the unit's own. Assets: no retained span contains a figure environment, a graphics-inclusion command, a PDF-page inclusion, a table or a TikZ picture; no image file is copied into this package, the package declares no asset dependency and no image file is redistributed with it. Licence: version arXiv:2410.17407v3 of this article is distributed under the Creative Commons Attribution 4.0 International licence, as recorded in the arXiv licence metadata for this exact version and shown on the abstract page https://arxiv.org/abs/2410.17407v3. A full rights notice is delivered at licenses/arxiv-2410.17407-CC-BY-4.0.txt. Characters: every retained excerpt is pure ASCII, so the delivered engine, XeLaTeX with its default Latin Modern text font, reports no missing character.
\begin{Prop}[Amplification]\label{amplification} Let $M,N$ be $W^*$-algebras. Then
$$M\sim N, \quad (A, \alpha)\sim_\G (B, \beta)\implies (M\ovot A, \id \otimes \alpha)\sim_\G (N \ovot B, \id \otimes \beta).$$
\end{Prop}

\begin{Theorem}\label{char1} 
    The following are equivalent:
    \begin{enumerate}\setlength\itemsep{-0.5em}
        \item There exists a $\G$-Morita correspondence $(\mathcal{H}, \pi, \rho, U)\in \Corr^\G(A,B)$.
        \item There exists a linking $\G$-$W^*$-algebra $(E, \gamma,p, \kappa, \lambda)$ between $(A, \alpha)$ and $(B, \beta)$.
    \end{enumerate}
\end{Theorem}

\begin{Exa}\label{main example}
    Let $(B, \beta)$ be a $\G$-$W^*$-algebra, $p\in B$ a full projection and $P\in B \ovot L^\infty(\G)$ satisfying
    $$(\id \otimes \Delta)(P) = (P\otimes 1) (\beta \otimes \id)(P), \quad P^*P=\beta(p), \quad PP^* = p\otimes 1.$$
    Then the following statements are true:
    \begin{enumerate}\setlength\itemsep{-0.5em}
        \item The assignment  $\beta_P: pBp \to pBp\ovot L^\infty(\G): x \mapsto P\beta(x)P^*$
    defines a $\G$-action $pBp \curvearrowleft \G$. 
    \item $(B, \beta)\sim_\G (pBp, \beta_P).$
    \end{enumerate}
\end{Exa}

\begin{Theorem}\label{char2} The following are equivalent:
      \begin{enumerate}\setlength\itemsep{-0.5em}
        \item $(A, \alpha)\sim_\G (B, \beta)$.
        \item There exists a Hilbert space $\mathcal{K}$, a full projection $p\in B(\mathcal{K})\ovot B$, $P\in B(\mathcal{K})\ovot B \ovot L^\infty(\G)$ and a $*$-isomorphism $\phi: p(B(\mathcal{K})\ovot B)p\to A$ such that the following conditions are satisfied:
    \begin{enumerate}\setlength\itemsep{-0.5em}
        \item $P^*P=(\id \otimes \beta)(p)$.
        \item $PP^* = p\otimes 1$.
        \item $(\id \otimes \id \otimes \Delta)(P) = (P\otimes 1)(\id \otimes \beta \otimes \id)(P).$
        \item $\alpha(\phi(z)) = (\phi\otimes \id)(P(\id \otimes \beta)(z)P^*)$ for all $z\in p(B(\mathcal{K})\ovot B)p.$
    \end{enumerate}
    \end{enumerate}
\end{Theorem}

\begin{proof} $(1)\implies (2)$ Let $(\mathcal{H}, \pi, \rho, U)\in \operatorname{Corr}^\G(A,B)$ be a $\G$-Morita correspondence. There exists a set $I$ and a projection $p\in B(\ell^2(I))\ovot B$ such that 
$$(\mathcal{H}, \rho)\cong ((\id \otimes \pi_B)(p)(\ell^2(I)\otimes L^2(B)), \tilde{\rho})$$
as right $B$-representations, where $\tilde{\rho}(b) = 1\otimes \rho_B(b)$ for $b\in B$.

Put $C:= B(\ell^2(I))\ovot B$, which we view as a $\G$-$W^*$-algebra for the amplified $\G$-action $\id \otimes \beta: C \to C \ovot L^\infty(\G)$. We prove now that $p\in C$ is full. Consider the two-sided $\sigma$-weakly closed ideal
$J = [CpC]^{\sigma\text{-weak}}.$
There is a central projection $r\in B$ such that $J = (1\otimes r)C$. We have
$$\tilde{\rho}(r^\perp)(\id \otimes \pi_B)(p)= (1\otimes \rho_B(r^\perp))(\id \otimes \pi_B)(p)= (\id \otimes \pi_B)((1\otimes r^\perp)p)= 0,$$
so that by faithfulness of $\tilde{\rho}$ we find that $r^\perp= 0$, so that $p$ is full in $C$.

Consider the $\G$-Morita-correspondence $\overline{\ell^2(I)}\otimes L^2(B)\in \Corr^\G(B,C)$ defined by
    \begin{itemize}\setlength\itemsep{-0.5em}
        \item the normal $*$-representation $\pi'(b) = 1_{\overline{\ell^2(I)}}\otimes \pi_B(b), \quad b \in B,$
        \item the normal anti-$*$-representation $\rho'(z) = ({(-)}^T \otimes \rho_B)(z), \quad z \in C,$
        \item the unitary $\G$-representation $U_{\beta, 23}\in B(\overline{\ell^2(I)}\otimes L^2(B))\ovot L^\infty(\G)$.
    \end{itemize}
We find
\begin{align*}
    \mathcal{H}\boxtimes_B (\overline{\ell^2(I)}\otimes L^2(B)) &\cong \overline{\ell^2(I)}\otimes \mathcal{H}\\
    &\cong \overline{\ell^2(I)}\otimes (\id \otimes \pi_B)(p)(\ell^2(I)\otimes L^2(B))\\
    &= (\pi_{B(\ell^2(I))}\otimes \pi_B)(p)(\overline{\ell^2(I)}\otimes \ell^2(I)\otimes L^2(B))\\
    &= \pi_{B(\ell^2(I))\ovot B}(p)L^2(B(\ell^2(I))\ovot B) = \pi_C(p)L^2(C)
\end{align*}
as right $C$-representations. Therefore, by structure transport, the Hilbert space $\mathcal{G}:=\pi_C(p)L^2(C)$ with its natural right $C$-action is endowed with the structure of a $\G$-$A$-$C$-Morita correspondence. The fact that $\mathcal{G}$ becomes a $\G$-$A$-$C$-Morita correspondence then implies that 
$$A \cong \pi_\mathcal{G}(A)= \rho_\mathcal{G}(C)' = (\pi_C(p) \pi_C(C) \pi_C(p))\vert_{\pi_C(p)L^2(C)}\cong \pi_C(pCp)\cong pCp,$$
so we find a $*$-isomorphism $\phi: pCp \to A$.

Consider the linking $\G$-$W^*$-algebra $E:=\mathscr{L}_C(\mathcal{G}\oplus L^2(C))$ between $(A, \alpha)$ and $(C, \id \otimes \beta)$ (see the proof of Theorem \ref{char1}) endowed with its natural $\G$-action $\gamma: E\to E\ovot L^\infty(\G)$
and the natural $\G$-equivariant embeddings
$$\kappa: A \to E: a \mapsto \begin{pmatrix}
    \pi_\mathcal{G}(a) & 0 \\ 0 & 0
\end{pmatrix}, \quad \lambda: C \to E: c \mapsto \begin{pmatrix}
    0 & 0 \\ 0 & \pi_C(c)
\end{pmatrix}.$$
Evidently, we can view
\begin{align*}
    \mathscr{L}_C(\mathcal{G}\oplus L^2(C))= \begin{pmatrix}
        \pi_C(pCp) & \pi_C(pC) \\ \pi_C(Cp) & \pi_C(C)
    \end{pmatrix},
\end{align*}
 where the right hand side acts naturally on $\mathcal{G}\oplus L^2(C)$. We note that 
$$ \gamma\begin{pmatrix}
    0 & \pi_C(p) \\ 0 & 0
\end{pmatrix}  = \begin{pmatrix}
    0 & (\pi_C\otimes \id)(P) \\ 0 & 0
\end{pmatrix}$$
for some element $P\in C \ovot L^\infty(\G)$ satisfying $P=(p\otimes 1)P$. Then note that 
$$\begin{pmatrix}
    0 & (\pi_C\otimes \id)(P) \\ 0 & 0
\end{pmatrix}^* \begin{pmatrix}
    0 & (\pi_C\otimes \id)(P) \\ 0 & 0
\end{pmatrix} = \gamma \begin{pmatrix}
    0 & 0 \\ 0 & \pi_C(p)
\end{pmatrix}= \begin{pmatrix}
    0 & 0 \\ 0 & (\pi_C \otimes \id)((\id \otimes \beta)(p))
\end{pmatrix}$$
implies the property $(a)$. Similarly, property $(b)$ is proven. Further,
\begin{align*}
    &\begin{pmatrix}
        0 & (\pi_C\otimes \id \otimes \id)((P\otimes 1)(\id \otimes \beta \otimes \id)(P))\\
        0 & 0
    \end{pmatrix}\\
    &= \begin{pmatrix}
        0 & (\pi_C\otimes \id)(P)\otimes 1\\ 0 & 0
    \end{pmatrix} \begin{pmatrix}
        0 & 0 \\ 0 &
         (\pi_C\otimes \id \otimes \id)(\id \otimes \beta \otimes \id)(P)
    \end{pmatrix}\\
    &= \left(\gamma\begin{pmatrix}
        0 & \pi_C(p)\\ 0 & 0
    \end{pmatrix}\otimes 1\right)(\gamma\otimes \id) \begin{pmatrix}
        0 & 0 \\ 0 & (\pi_C\otimes \id)(P)
    \end{pmatrix}\\
    &= (\gamma \otimes \id)\begin{pmatrix}
        0 & (\pi_C\otimes \id)(P) \\ 0 & 0
    \end{pmatrix}\\
    &= (\id \otimes \Delta)\begin{pmatrix}
        0 & (\pi_C\otimes \id)(P) \\ 0 & 0
    \end{pmatrix}= \begin{pmatrix}
        0 & (\pi_C\otimes \id \otimes \id)(\id_C \otimes \Delta)(P)\\
        0 & 0
    \end{pmatrix},
\end{align*}
from which $(c)$ follows. Next, note that for $c\in C$, 
\begin{align*}
    (\kappa \otimes \id)\alpha(\phi(pcp))&=  \gamma\begin{pmatrix}
        \pi_C(pcp) & 0 \\ 0 & 0
    \end{pmatrix}\\
    &=\gamma\left(\begin{pmatrix}
        0 & \pi_C(p) \\ 0 & 0
    \end{pmatrix} \begin{pmatrix}
        0 & 0 \\ 0 & \pi_C(c)
    \end{pmatrix} \begin{pmatrix}
        0 & 0 \\ \pi_C(p) & 0
    \end{pmatrix}\right)\\
    &= \begin{pmatrix}
        0 & (\pi_C\otimes \id)(P)\\ 0 & 0
    \end{pmatrix} \begin{pmatrix}
        0 & 0 \\ 0 & (\pi_C \otimes \id)((\id \otimes \beta)(c))
    \end{pmatrix} \begin{pmatrix}
        0 & 0 \\ (\pi_C\otimes \id)(P^*) & 0
    \end{pmatrix}\\
    &= \begin{pmatrix}
        (\pi_C\otimes \id)(P(\id \otimes \beta)(c)P^*) & 0 \\
        0 & 0
    \end{pmatrix}\\
    &= (\kappa \otimes \id)(\phi\otimes \id)(P(\id \otimes \beta)(c)P^*),
\end{align*}
from which $(d)$ follows. 

$(2)\implies (1)$ We have
$$(B, \beta) \sim_\G (B(\mathcal{K})\ovot B, \id \otimes \beta) \sim_\G (p(B(\mathcal{K})\ovot B)p, (\id \otimes \beta)_P) \stackrel{\phi}{\cong}_\G (A, \alpha),$$
where the first $\G$-$W^*$-Morita equivalence follows from Proposition \ref{amplification} and the second $\G$-$W^*$-Morita equivalence follows from Example \ref{main example}.
\end{proof}

\end{document}
